% PATCH: R3-03
% Reviewer: 3
% Target section: sections/methodology.tex
% Requirement: "The low-torque servo constraints and open-loop gait defects are only briefly
%   mentioned, without supplementary improved control strategies or compensating algorithms to
%   address the platform's [limitations]." — results.tex's Limitations subsection already
%   acknowledged the constraint but proposed no compensation strategy; this patch adds the
%   strategy itself as a new subsection at the end of methodology.tex.
% Status: applied
%
% Source material: reviewed Peter_arduino/src firmware (gait.cpp, imu.cpp, kinematics.cpp,
% config.h, peter_arduino.ino) from the PETER_SIMULATION reference repo (intake/SOURCES.md).
% Core three techniques included (independent of any pending confirmation from the firmware
% developer): (1) statically-stable single-leg-swing wave gait -- structural CoM/support-triangle
% guarantee, cross-referenced to the existing Gait Decision Module description; (2) rate-limited
% (velocity-capped) trajectory interpolation -- bounds peak torque/current demand; (3) fixed
% feed-forward tibia-angle offset compensating measured servo sag under load. Four supporting
% measures (per-actuator PWM calibration, software angle saturation, IK-feasibility rejection)
% mentioned briefly in one sentence as lower-level safeguards, not as primary findings.
%
% Deliberately EXCLUDED pending confirmation from the developer (Sam): the IMU-based tilt safety
% stop and its complementary-filter attitude estimation/bias-nulling support -- config.h ships
% with USE_IMU=0 by default ("sin correccion ni parada de seguridad"), so it cannot yet be claimed
% these operated during the reported real-robot experiments. Also excluded: IMU_GAIN_X/Y (defined
% in config.h but never referenced anywhere in the codebase -- designed, not implemented).
%
% Closing sentence explicitly notes this does not substitute for closed-loop torque/attitude
% feedback during active balance recovery, cross-referencing the existing Limitations subsection
% (Section~\ref{ssec:Limitations}) -- keeps the new content consistent with, not contradicting,
% that subsection's existing "open-loop execution cannot compensate for torque deficits [on severe
% terrain]" statement (a narrower claim about active recovery on steep/high-slip terrain, not a
% blanket claim of zero compensation anywhere).

\subsection{\revblue{Low-Torque Actuation Compensation Strategies}}
\label{ssec:ServoCompensation}

\revblue{The MG996R servomotors driving the leg joints (Section~\ref{ssec:Construction}, nominal torque $11~\text{kg·cm}$) are commanded in open loop, with no torque or current feedback at the joint level. Under the sustained load of quadrupedal support and locomotion, this combination is prone to mechanical compliance ---a deviation between the commanded joint angle and the angle physically achieved--- and to abrupt-command-induced instability if the trajectory generator demands more torque than the actuator can deliver instantaneously. Three complementary strategies, implemented directly in the gait-execution firmware, mitigate these effects without requiring closed-loop torque control.}

\revblue{First, gait stability is addressed at the scheduling level rather than through active balance control. The crawl gait described in Section~\ref{ssec:Gait_Decision_Module} advances exactly one leg at a time, following a fixed diagonal sequence (front-left~$\rightarrow$~rear-right~$\rightarrow$~front-right~$\rightarrow$~rear-left); the remaining three legs stay grounded and fixed throughout each swing phase. This is a structural, not a runtime-computed, guarantee: because only one leg is ever airborne, the robot's center of mass is always contained within the support triangle formed by the three stance legs, independently of any per-step positioning error the open-loop actuation may introduce. The gait is therefore quasi-statically stable by construction, which bounds the destabilizing consequence of an individual joint's torque deficit to a single, momentarily unsupported limb rather than to the whole platform.}

\revblue{Second, foot trajectories are interpolated toward each target position at a velocity capped per control iteration (rather than commanded as an instantaneous step change), separately for the swing-leg motion and for the whole-body stride-back phase. Bounding the commanded rate of change caps the peak torque and current the actuator must instantaneously supply to track the trajectory, reducing the risk of position lag or oscillation that an abrupt setpoint jump would otherwise induce on a low-torque actuator.}

\revblue{Third, a fixed corrective offset is added to the commanded tibia angle before it is sent to the actuator, empirically determined to counteract the servo's mechanical sag under the knee joint's load: the achieved foot position is thereby brought back into agreement with the kinematic target computed by the inverse-kinematics solver, rather than settling short of it as the uncompensated command would. This is a static, feed-forward correction, applied uniformly at every control cycle rather than derived from a real-time load or position measurement.}

\revblue{These three measures are supplemented by lower-level safeguards common to servo-actuated platforms: an individually calibrated pulse-width mapping per actuator (compensating unit-to-unit mechanical and mounting variation across the twelve joints), a software saturation limit on every commanded joint angle, and a rejection of inverse-kinematics solutions that fall outside the physically reachable workspace before any command reaches the actuator. Together, this layered strategy --- static stability by gait design, torque-aware trajectory shaping, and feed-forward compliance correction --- constitutes the compensating control approach for the platform's low-torque, open-loop actuation constraints. It does not substitute for closed-loop torque or attitude feedback during active balance recovery on severe terrain, which remains an open limitation of the present design (Section~\ref{ssec:Limitations}).}
